\subsection{Background and Motivation}

Robotic manipulators are increasingly deployed in domains such as industrial automation, teleoperation, and robot-assisted surgery. In many of these applications, operator training is essential to ensure safe and efficient task execution. While visual feedback provides global situational awareness, it is often insufficient for developing fine motor skills, accurate force regulation, and intuitive understanding of interaction constraints.

Haptic interfaces enable bidirectional physical interaction between a human operator and a robotic or simulated environment. By rendering interaction forces to the user, such systems can improve perception of contact conditions, enhance motor learning, and support safer manipulation strategies during training tasks. As robotic systems become more widely adopted in both research and educational settings, there is increasing motivation to develop accessible haptic training platforms that allow users to experience realistic interaction dynamics without requiring direct access to expensive physical manipulators.

Despite these motivations, many existing haptic training solutions remain costly, proprietary, and tightly coupled to specific software ecosystems. This limits transparency of the control architecture and restricts the ability of researchers and students to experiment with alternative control strategies, custom simulation environments, or novel hardware configurations. There is therefore a practical need for modular and low-cost haptic interface systems that enable real-time bidirectional interaction with robotic manipulation simulations while remaining flexible and extensible.

% NOTE: keep this section conceptual and problem-driven.
% Detailed discussion of prior work and specific rendering methods
% has been intentionally moved to the literature review.

\subsection{Research Gap}

While prior work has addressed either high-fidelity commercial haptic interfaces or low-cost robotic hardware in isolation, there remains limited work integrating low-cost actuated haptic devices with transparent and modular software architectures suitable for rapid prototyping and educational use.

In particular, there is a practical need for systems that combine:
\begin{enumerate}
	\item Low-cost hardware, with total prototype costs significantly below commercial systems that typically exceed £5{,}000--£20{,}000.
	\item A modular communication and software architecture capable of supporting integration with multiple simulation backends.
	\item Explicit control over latency, stability, and force rendering behaviour.
	\item A software framework that can be extended toward physical robotic manipulators in future work.
\end{enumerate}

\subsection{Aims and Objectives}

The aim of this project is to design, implement, and evaluate a low-cost, open-architecture two-degree-of-freedom (2-DOF) haptic interface and simulation framework capable of delivering real-time, stable force feedback for robotic manipulator training.

The specific objectives are to:
\begin{itemize}
	\item Design and construct a 2-DOF haptic input device using low-cost BLDC actuators and magnetic rotary encoders.
	\item Develop embedded firmware implementing field-oriented motor control, encoder acquisition, and real-time bidirectional communication.
	\item Define and implement a device-agnostic packet-based communication protocol for streaming joint states and receiving force or torque commands.
	\item Develop a custom host-side real-time software framework built around NVIDIA PhysX and OpenGL, supporting collision detection, virtual proxy methods, and real-time force computation.
	\item Implement Jacobian-based mappings from end-effector interaction forces to joint-space actuator torques.
	\item Evaluate system performance in terms of latency, stability, and force fidelity against established haptic rendering requirements.
	\item Demonstrate an architecture designed to support future integration with external simulation environments such as NVIDIA Isaac Sim or MATLAB/Simulink.
\end{itemize}

\subsection{Contributions of This Work}

This work presents the design and implementation of a modular real-time haptic rendering system integrating a low-cost actuated input device with a custom multi-threaded host-side software framework.

The complete hardware prototype was realised at an estimated cost of approximately £180, substantially below the cost of commercial haptic devices such as the Geomagic Touch, Touch X, and Force Dimension systems, which typically range from several thousand to tens of thousands of pounds.

The software architecture was designed to be simulation-backend-independent through the use of a device-agnostic messaging interface and modular subsystem separation, enabling future integration with external platforms such as MATLAB/Simulink or NVIDIA Isaac Sim.

The principal contributions of this work are:

\begin{itemize}
	\item Design and construction of a low-cost 2-DOF actuated haptic interface.
	\item Development of a custom message-passing architecture for deterministic inter-thread communication.
	\item Development of a custom proxy-based haptic rendering engine using signed distance field collision queries.
	\item Integration of rigid body simulation using NVIDIA PhysX through a custom wrapper layer.
	\item Development of a custom OpenGL-based rendering and visualisation framework.
	\item Explicit consideration of latency, stability, and safety mechanisms for real-time haptic interaction.
\end{itemize}

% NOTE: refine contributions AFTER results section is complete.
% Add quantitative achievements here later (latency, force bandwidth, etc).

% NOTE: also add a short "Report Structure" paragraph at the end of the intro.